\section{PAPER 029: ProvenanceLedger — Targeted Rollback and Audit for Autonomous Knowledge Graph Materi\-alization}

\textbf{Author}: Bryan Alexander Buchorn  
\textbf{Affiliation}: Independent Researcher, Las Vegas, NV, USA  
\textbf{Status}: v2.51.0 (Phase 167 (STRB) COMPLETE)
\textbf{Date}: May 2, 2026

\textbf{CEREBRUM Phase 76}

---

\subsection{Abstract}

We present \textbf{ProvenanceLedger}, an audit chain component for CEREBRUM's autonomous knowledge graph discovery pipeline. ProvenanceLedger records every edge mater\-ialized by \texttt{ResearchAgent.approve()} in structured \texttt{EdgeRecord} / \texttt{BatchRecord} objects, enabling targeted rollback at two granu\-larities: per-approval-batch (\texttt{rollback\_batch}) and per-loop-cycle (\texttt{rollback\_cycle}). Unlike naive graph snapshots, which require full graph comparison to identify removed edges, ProvenanceLedger provides O(1) lookup of edges-to-remove by batch ID or cycle number. The ledger is thread-safe, LRU-capped at \texttt{max\_batches} to bound memory usage, and requires adapters to implement the Phase 80 \texttt{remove\_edge()} protocol. ProvenanceLedger is the prere\-quisite for Loop-Provenance Recovery (Phase 79), which wires circuit breaker trips to automatic cycle rollback.

---

\subsection{1. Motivation: Accoun\-tability in Autonomous Materi\-alization}

Prior to Phase 76, \texttt{ResearchAgent.approve()} called \texttt{adapter.add\_edge()} with no record of which edges were added in which approval decision. If a batch of findings later proved incorrect (e.g., when the circuit breaker tripped, indicating poor discovery quality), there was no automated way to remove exactly those edges without manual graph inspection.

Two rollback granu\-larities are needed:
\begin{enumerate}
\item \textbf{Batch rollback}: Remove edges from one specific \texttt{approve()} call. Useful when a single finding is later disproved.
\item \textbf{Cycle rollback}: Remove all edges mater\-ialized in loop cycle N. Useful when the entire cycle is suspect (circuit breaker trip, external validation failure).

\end{enumerate}
---

\subsection{2. Data Model}

\subsubsection{EdgeRecord}
\begin{lstlisting}
@dataclass
class EdgeRecord:
    u: str           # source entity
    v: str           # target entity
    relation: str    # edge label
    finding_id: str  # originating ResearchFinding ID
\end{lstlisting}

\subsubsection{BatchRecord}
\begin{lstlisting}
@dataclass
class BatchRecord:
    batch_id: str              # timestamp + finding hash
    cycle_number: int          # loop cycle that produced this batch
    edges: List[EdgeRecord]
    rolled_back: bool = False
    rolled_back_at: Optional[float] = None
\end{lstlisting}

\subsubsection{ProvenanceLedger}
\begin{itemize}
\item \texttt{\_batches: OrderedDict[str, BatchRecord]} — LRU-ordered dict
\item \texttt{\_cycle\_index: Dict[int, List[str]]} — cycle $\to$ list of batch\_ids

\end{itemize}
---

\subsection{3. Recording}

\begin{lstlisting}
# Called automatically by ResearchAgent.approve()
ledger.record(batch_id, cycle_number, edges: List[Tuple[str, str, str]])
\end{lstlisting}

When \texttt{max\_batches} is reached, the oldest batch is evicted (LRU).

---

\subsection{4. Rollback}

\subsubsection{Batch rollback}
\begin{lstlisting}
removed: int = ledger.rollback_batch("batch_20260414_001", adapter)
# Calls adapter.remove_edge(u, v, relation) for each EdgeRecord
# Marks batch as rolled_back=True
\end{lstlisting}

\subsubsection{Cycle rollback}
\begin{lstlisting}
removed: int = ledger.rollback_cycle(12, adapter)
# Iterates all batches with cycle_number == 12
# Calls rollback_batch for each
\end{lstlisting}

Both methods return the count of edges removed, which is recorded in \texttt{CycleRecord.edges\_rolled\_back} (Phase 79).

---

\subsection{5. Stats and Inspection}

\begin{lstlisting}
stats = ledger.stats()
# {
#   "total_batches": 42,
#   "active_batches": 38,
#   "rolled_back_batches": 4,
#   "total_edges": 187,
#   "rolled_back_edges": 19,
# }

batches = ledger.recent_batches(n=10)
# List of BatchRecord dicts, newest first
\end{lstlisting}

---

\subsection{6. REST API}

\begin{center}
{\footnotesize
\begin{tabularx}{\columnwidth}{|>{\raggedright\arraybackslash}X|>{\raggedright\arraybackslash}X|>{\raggedright\arraybackslash}X|}
\hline
Endpoint & Method & Description \\ \hline
\texttt{/research/provenance/stats} & GET & Ledger totals \\ \hline
\texttt{/research/provenance/batches} & GET & Recent batches list (\texttt{?n=20}) \\ \hline
\texttt{/research/provenance/rollback/{batch\_id}} & POST & Rollback one batch \\ \hline
\texttt{/research/provenance/rollback-cycle/{n}} & POST & Rollback all batches from cycle N \\ \hline
\end{tabularx}
}
\end{center}

---

\subsection{7. Thread Safety}

All read and write operations on \texttt{\_batches} and \texttt{\_cycle\_index} are protected by a \texttt{threading.Lock}. Concurrent \texttt{approve()} calls from different threads produce distinct \texttt{batch\_id} values (based on \texttt{time.time\_ns()} + \texttt{uuid4} suffix) and never collide.

---
\textbf{Copyright © 2026 Bryan Alexander Buchorn. All Rights Reserved.}

---
\subsection{Acknow\-ledgments}

The author gratefully ackno\-wledges the use of Claude (Anthropic) as a research assistant throughout this work. Claude assisted with mathe\-matical forma\-lization, code generation, manuscript preparation, and technical writing. All conceptual contr\-ibutions, archi\-tectural decisions, exper\-imental design, and intel\-lectual claims are solely the author's.

\textbf{Reviewed on}: May 2, 2026 for version v2.51.0
